% !PRINTED 71
\weberpage{71}
%<<CONT
erhält. Das ist aber die Gruppe
\weq{7}{\mathfrak{A}_\mu \equiv 1\,(\mathrm{mod}\ \mu),}

d. h. die Gruppe, die aus allen den Zahlen von $\mathfrak{N}$ besteht, die
der Congruenz $a \equiv 1\,(\mathrm{mod}\ \mu)$ genügen. Denn zunächst ist klar,
dass $\mathfrak{A}_{m_1}$ und $\mathfrak{A}_{m_2}$ Divisoren von $\mathfrak{A}_\mu$ sind. Also ist auch das
kleinste gemeinschaftliche Vielfache von $\mathfrak{A}_{m_1}$ und $\mathfrak{A}_{m_2}$ in $\mathfrak{A}_\mu$
enthalten. Ist aber andererseits $\alpha = 1 + \xi\mu$ irgend eine Zahl in $\mathfrak{A}_\mu$, so
kann man die Zahlen $a_1 = 1 + x_1 m_1$ in $\mathfrak{A}_{m_1}$ und $a_2 = 1 + x_2 m_2$
in $\mathfrak{A}_{m_2}$ so bestimmen, dass $\alpha \equiv a_1 a_2\,(\mathrm{mod}\ m)$ wird. Man hat
nur $\xi\mu = x_1 m_1 + x_2 m_2$ zu setzen, was nach Bd. I, §. 118 immer
möglich ist. Es ist also auch $\mathfrak{A}_\mu$ in dem kleinsten gemeinschaftlichen Vielfachen von $\mathfrak{A}_{m_1}$ und $\mathfrak{A}_{m_2}$ enthalten, welches demnach
$\mathfrak{A}_\mu$ selbst ist. Daraus also der Satz:

\gesperrt{4. Sind $m_1$ und $m_2$ irgend zwei natürliche Zahlen
und $\mu$ der grösste gemeinschaftliche Theiler von
$m_1$ und $m_2$, so ist der Körper $\Omega_\mu$ der grösste
gemeinschaftliche Theiler von $\Omega_{m_1}$ und $\Omega_{m_2}$ und
$\Omega_{\frac{m_1 m_2}{\mu}}$ ihr kleinstes gemeinsames Multiplum.}

Daraus folgt nun: wenn zwei volle Kreistheilungskörper $\Omega_m$
und $\Omega_{m'}$ einen gemeinsamen Theiler $\Omega$ haben, und $m'$ ist kleiner
als $m$, so muss es einen echten Theiler $d$ von $m$ geben, so dass
$\Omega$ auch ein Theiler von $\Omega_d$ ist.

Wir wollen einen Theiler von $\Omega_m$ \gesperrt{primär} nennen, wenn er
nicht zugleich in einem vollen Kreistheilungskörper $\Omega_{m'}$ von
niedrigerem $m'$ enthalten ist. Dann folgt also, dass man alle
nicht primären Theiler von $\Omega_m$ erhält, wenn man in $\Omega_d$ den
Index $d$ alle echten Theiler von $m$ durchlaufen lässt und die
Theiler von $\Omega_d$ aufsucht.

Bezeichnen wir nun mit $q_1, q_2, \ldots q_\tau$ die sämmtlichen in $m$
aufgehenden verschiedenen Primzahlen (2 eingeschlossen) und
setzen
\weq{8}{m = q_1 m_1 = q_2 m_2 = \cdots = q_\tau m_\tau,}

so ist jedes $d$ Theiler von einem der $m_1, m_2, \ldots m_\tau$, und wir erhalten die nicht primären Theiler von $\Omega_m$, wenn wir alle Theiler
der Körper
\weq{9}{\Omega_{m_1},\ \Omega_{m_2},\ \ldots\ \Omega_{m_\tau}}

aufsuchen. Diese Körper gehören aber zu der Gruppe
\weq{10}{Q_1 \equiv 1\,(\mathrm{mod}\ m_1),\ \ Q_2 \equiv 1\,(\mathrm{mod}\ m_2),\ \ldots\ Q_\tau \equiv 1\,(\mathrm{mod}\ m_\tau),}
%CONT>>
